% SEGMENT OLP-0659-S01
% Part: proof-theory
% Chapter: cut-elimination
% Section: introduction

\documentclass[../../../include/open-logic-section]{subfiles}

\begin{document}

\olfileid{pt}{cut}{int}

\olsection{அறிமுகம்}

\Cut{} விதி பின்வருவது:
\[
\Axiom$\Gamma \fCenter \Delta, !A$
\Axiom$!A, \Pi \fCenter \Lambda$
\RightLabel{\Cut}
\BinaryInf$\Gamma, \Pi \fCenter \Delta, \Lambda$
\DisplayProof
\]
(இங்குள்ள !!{formula}~$!A$ \,\emph{வெட்டு வாய்பாடு} எனப்படுகிறது.)

கென்ட்சனின் முதல் தொடரணிக் கணியமான~\Log{LK} இல்
\Cut{} சேர்க்கப்பட்டிருந்தது. கென்ட்சனின்
\emph{Hauptsatz} (``முதன்மைத் தேற்றம்'') கூறுவது:
\Log{LK} இல் !!{provable} ஆன ஒவ்வொரு தொடரணியும் \Cut{} விதியைப்
பயன்படுத்தாமலும் !!{provable} ஆகும். அதாவது, \Log{LK}-நிறுவல்களிலிருந்து
\Cut{} ஐ ``நீக்க'' முடியும். நமது~\Log{G1c}, \Log{G3c}
முறைமைகளில் \Cut{} விதி இல்லை. இருந்தாலும் அவற்றுக்கும் அதே கேள்வியை
எழுப்பலாம்: \Log{G1c} (அல்லது \Log{G3c}) இன் மற்ற விதிகளுடன்
\Cut{} ஐயும் பயன்படுத்தும் !!{proof}s இலிருந்து \Cut{} ஐ
நீக்க முடியுமா? ஏற்கெனவே வகுத்த சொல்லாட்சியில், இது
\Cut{} ஆனது \Log{G1c} (அல்லது \Log{G3c}) இல்
ஏற்கத்தக்க விதியா என்று கேட்பதற்குச் சமம்.

% SEGMENT OLP-0659-S02
வெட்டு நீக்கத் தேற்றம் நிறுவல் கோட்பாட்டின் மிக முக்கியமான,
பரவலாகப் போற்றப்படும் முடிவுகளில் ஒன்று. முதல் பயனாக,
\Log{G1c}, \Log{G3c} ஆகியவற்றிற்கு, முதலில் பார்த்தால்
இன்றியமையாததாகத் தோன்றும் ஒரு விதி உண்மையில் தேவையில்லை என்று
இது காட்டுகிறது. உண்மையான கணித நிறுவல்களை முறையாக்கும்போது
துணைத்தேற்றங்களைப் பயன்படுத்துவது எப்படி என்ற கேள்வி எழும்.
கணிதவியலாளர்கள் ஒரு முடிவை நிறுவி, அடுத்த முடிவின் நிறுவலில்
அதைப் பயன்படுத்துகிறார்கள். முதலில் துணைத்தேற்றமான~$L$ இன்
நிறுவலை $\Sequent L$ என்ற தொடரணியின் !!a{proof} ஆக
முறையாக்கலாம். பிறகு~$L$ ஐப் பயன்படுத்தி இரண்டாம் முடிவான~$R$ ஐ
நிறுவுவதை $L \Sequent R$ என்ற தொடரணியின் !!a{proof} ஆக
முறையாக்கலாம். அப்படியானால், $R$ க்கே, அதாவது $\Sequent R$ க்கே,
!!a{proof} இருப்பதை எப்படித் தெரிந்துகொள்வது? இரண்டு
!!{proof}s ஐ இணைக்க வழி வேண்டும்; அதை \Cut{} விதி தருகிறது.

% SEGMENT OLP-0659-S03
\Log{G1c}, \Log{G3c} ஆகியவை நிறைவானவை என முன்பு
குறிப்பிட்டோம். எனவே நிறுவத்தக்கதாக எதிர்பார்க்கும் ஒவ்வொரு
தொடரணியையும், அதாவது செல்லுபடியான எல்லாத் தொடரணிகளையும்,
இம்முறைமைகள் நிறுவுகின்றன. இதை நேரடியாக நிறுவலாம். ஆனால்
கென்ட்சன் எழுதிய காலத்தில், அடிக்கோள் சார்ந்த வருவித்தல்
முறைமைகளின் நிறைவுமை மட்டுமே அறியப்பட்டிருந்தது.
\Log{LK} இன் நிறைவுமையைக் காட்ட, அடிக்கோள் சார்ந்த முறைமையின்
!!{derivation}s ஐ \Log{LK}-!!{proof}s ஆக மாற்றும் முறையை
கென்ட்சன் தந்தார். அம்மாற்றத்தில் \Cut{} இன்றியமையாததாக
இருந்தது. வெட்டு நீக்கம் செவ்வியல் தருக்கத்திற்கு மட்டும்
முக்கியமானதன்று: வேறு பல தருக்கங்களுக்கும் தொடரணிக் கணியங்கள்
உள்ளன. சில தருக்கங்களில் அக்கணியங்களின் நிறைவுமையை நேரடியாக
நிறுவுவது கடினமாகவோ இயலாததாகவோ இருக்கும். அப்போது \Cut{} ஐப்
பயன்படுத்தி வேறு முறைமையிலிருந்து மாற்றுவதுதான் நிறைவுமையைக் காட்டும்
ஒரே வழியாக இருக்கலாம். பின்னர், \Cut{} தேவையற்றது என்றும்
\Cut{} இல்லாத முறைமையும் நிறைவானது என்றும் நிறுவ,
நிறுவல்-கோட்பாட்டு முறையில் வெட்டு நீக்கத்தை நிறுவ வேண்டும்.

% SEGMENT OLP-0659-S04
வெட்டு நீக்கத்தின் மற்றொரு முக்கியத்துவம், அதிலிருந்து தம்மளவில்
சுவாரசியமான பிற முடிவுகளையும் பெறலாம் என்பதே. நாம் பார்க்கவிருக்கும்
இரண்டு எடுத்துக்காட்டுகள் கிரெய்கின் இடையீட்டுத் துணைத்தேற்றமும்
ஹெர்பிராண்டின் தேற்றமும் ஆகும்.

வெட்டு நீக்கத்தை நிறுவும் கருத்து பொதுவாக, \Cut{} ஐப் பயன்படுத்தும்
!!a{proof} ஐ அது இல்லாத நிறுவலாக மாற்றுவது எப்படி என்பதைக்
காட்டுவதாகும். இம்மாற்றங்கள் பெரும்பாலும் ``பகுதிசார்'':
ஒரு நிறுவலின் ஒரு பகுதியை, பொதுவாக \Cut{} அனுமானத்தில் முடியும்
உள்-!!{proof} ஐ, அதே தொடரணியை நிறுவும் வேறோர்
உள்-!!{proof} ஆல் மாற்றுகிறோம். புதிய உள்-நிறுவலில் அந்த
\Cut{} அனுமானம் நீக்கப்பட்டிருக்கும் அல்லது வேறு அனுமானங்களால்
மாற்றப்பட்டிருக்கும். இவ்வகை வாதங்கள், \Cut{} உள்ள
!!{proof}s ஐ வெட்டு இல்லாதவற்றாகப் படிப்படியாக மாற்றும்
படிமுறைகளையும் தருகின்றன. பயன்பாடுகளில் வெட்டு நீக்க நடைமுறைகளைப்
பயன்படுத்தும்போது இது மேலும் தகவல் தருகிறது. எடுத்துக்காட்டாக,
$!A \lif !B$ செல்லுபடியானது எனில் ஓர் இடையீட்டி இருக்கிறது
என்ற வடிவில் இடையீட்டுத் துணைத்தேற்றத்திற்கு மாதிரிக் கோட்பாட்டு
நிறுவல்கள் உள்ளன (இடையீட்டி என்றால் என்ன என்பதை இப்போது
ஒதுக்கிவைப்போம்). வெட்டு நீக்கத்தைப் பயன்படுத்தும் நிறுவல்-கோட்பாட்டு
வாதமோ $!A \lif !B$ இன் !!a{proof} ஒன்றிலிருந்து இடையீட்டியைக்
கணக்கிட்டுத் தரும் படிமுறையை வழங்குகிறது. மாதிரிக் கோட்பாட்டு
வாதத்துடன் ஒப்பிடும்போது, இந்த நிறுவல்-கோட்பாட்டு வாதம்
\emph{கட்டுமானத் தன்மை உடையது}.

% SEGMENT OLP-0659-S05
வெட்டு நீக்கத் தேற்றத்தை நிறுவ இரண்டு பரவலான முறைகள் உள்ளன.
கென்ட்சனிடமிருந்து வந்த முதல் முறை, ஒரு நிறுவலில் உச்சியிலுள்ள
வெட்டுகளை நீக்க முடியும் எனக் காட்டி, எதுவும் எஞ்சாதவரை அவற்றை
ஒன்றன்பின் ஒன்றாக நீக்குகிறது. முதலில் ஷூட்டெவும் டைட்டும் வழங்கிய
இரண்டாம் முறை, !!a{proof} ஒன்றின் மிகச் சிக்கலான வெட்டுகளில் கவனம்
செலுத்துகிறது. அதன் வெட்டு வாய்பாட்டை விட அதிக ஆழமுள்ள
வெட்டு !!{formula} வேறு எந்த வெட்டுக்கும் இல்லாதபோது அதை மிக உயர்ந்த
தரமுள்ள வெட்டாகக் கொண்டு, அத்தகைய வெட்டுகளைத் தொடர்ந்து
நீக்குகிறது. இரு முறைகளுக்கும் தனித்தனி நன்மைகளும் குறைகளும் உள்ளன.
சில முறைமைகளில் ஒரு முறை இயங்க, மற்றொன்றைச் செயல்படுத்துவது
கடினமாகவோ இயலாததாகவோ இருக்கலாம். ஆகவே இரண்டு முறைகளையும்
விளக்குகிறோம். \Log{G3c} க்கு வெட்டு நீக்கத்தை நிறுவுவோம்.

% SEGMENT OLP-0659-S06
உண்மையில், நாம் நீக்கக்கூடியது \Cut{} என்பதை விட,
\emph{சூழலைப் பகிரும் வெட்டு}~\CutCS{} என்பதை நிறுவுவோம்:
\[
\Axiom$\Gamma \fCenter \Delta, !A$
\Axiom$!A, \Gamma \fCenter \Delta$
\RightLabel{\Cut}
\BinaryInf$\Gamma \fCenter \Delta$
\DisplayProof
\]
ஒரு தொடரணிக் கணியம் வலுவிழக்கச் சேர்த்தலையும் சுருக்கத்தையும்
அனுமதித்தால், அவை \Log{G1c} இல் நேரடி விதிகளாகவோ
\Log{G3c} இல் ஏற்கத்தக்க விதிகளாகவோ இருக்கலாம்; அப்போது
\Cut{} கொண்டு \CutCS{} ஐச் செயல்படுத்தலாம்; மறுதிசையிலும்
அவ்வாறே செய்யலாம். முதல் வெட்டு நீக்க முறையை \CutCS{} க்கு
விளக்குவது எளிது; இரண்டாம் முறை \CutCS{} க்கே பொருந்தும்.
ஆகவே \Cut{} ஐ விட \CutCS{} ஐத் தேர்ந்தெடுக்கிறோம்.

\end{document}
